% Physical page 5 onward: the developed commentary of the concise study,
% organized by the questions the three interpretations answer differently
% rather than by the order of the Mass. Each unit draws on several appointed
% elements and sets the interpretations' answers beside one another.

\section{The Propers: Detailed Commentary}
\zlabel{triptych:concise:commentary:start}

\subsection{What is owed to God? The coin, the image and the whole self}

``To God, the things that are God's'' names no object, and the questioners had
not asked about it. The witnesses supply the object in different ways.

St Jerome gives the plain sense: ``Render to Caesar the things that are
Caesar's, that is, coin, tribute and money; and to God the things that are
God's: let us understand tithes, first-fruits, oblations and victims''
(\work{Commentarii in Matthaeum} III). The saying sets the dues of the temple
beside the dues of the empire, and the image belongs to the coin alone. St
Augustine takes another step, and in his exposition of Psalm~4 shows where it
comes from. The light of God's countenance is stamped on us ``as a penny is
stamped with the king's image. For man was made after the image and likeness of
God, which he defaced by sin: therefore it is his true and eternal good, if by
a new birth he be stamped.'' Then: ``I believe this to be the bearing of that
which some understand skilfully'', that ``as the tribute money is rendered to
him, so should the soul to God, illumined and stamped with the light of His
countenance'' (\work{Enarrationes in Psalmos} 4.8). The inference is offered
as what ``some'' understand, it rests on Genesis 1:26 and not on anything Jesus
says, and it looks past the image as given to an image defaced and stamped
again. In the tractates on John it becomes a story: ``We are God's money: we
have wandered away as coin from the treasury. The impression that was stamped
upon us has been rubbed out by our wandering. He has come to refashion'': ``to
Caesar his money, to God yourselves'' (40.9). On ``I will pay my vows to the
Lord'' Augustine draws the command, ``let him vow himself, let him pay himself
as a vow: this is exacted, this is due'' (\work{Enarrationes in Psalmos}
115.13).

St Ambrose, on Luke's form of the scene, moves the inference from the believer
to its source. ``The image of God is one thing, the image of the world
another'', and Paul bids us
bear ``the image of the heavenly'' (1~Cor 15:49); \latin{Imaginem Caesaris non
habet Christus, quia imago est Dei}: Christ does not have Caesar's image,
because he is the image of God, and when he paid the census he ``did not give
of his own, but returned to the world what was the world's''
(\work{Expositio evangelii secundum Lucam} IX.34--35). St Hilary of Poitiers
reaches the same object by another road. What is properly God's is
\latin{corpus, animam, voluntatem}, body, soul and will, and his reason is not
the image but origin: men are to render themselves whole to the one to whom
they owe ``both their origin and their growth'' (\work{Commentarius in
Matthaeum} 23.2). Augustine's hinge is the coin's image, Hilary's that the self
came from God; the two arguments meet in the whole self from different
starting points, and Hilary's list names what Augustine's sentences do not, the
will. St Thomas Aquinas holds a literal and a mystical sense together: natural
goods such as bread and wine are given to God and coin to Caesar, and
\latin{Mystice sic. Nos habemus animam quae est ad imaginem Dei, ideo eam Deo
reddere debemus}, mystically, we owe God the soul made to his image
(\work{Super Matthaeum} c.~22). In Aquinas the two readings do not compete.

Jerome's literal answer ends close to the others. Christ rendered to God
\latin{Patris faciens voluntatem}, doing the will of the Father, and on that
point the literal reader and Hilary meet at the will without the image. Cyril
of Alexandria says the same of Luke's scene: God ``requires of us of things
corruptible and transitory even nothing, but rather willing obedience and
submission; faith and love'' (\work{Commentary on Luke}, Sermon 135). The first
interpretation carries Augustine's inference through the whole Mass: Caesar's
claim reaches the coin, God's claim reaches the one who carries it, and the
Collect's will devoted to God becomes the measure of the Gospel's last words.
Its strain is the distance between the saying and the inference, for Jesus
speaks of a coin and of God, not of man as God's coin. Aquinas's two senses
show the reply: the literal sense stands beside the mystical one, and leaves
``the things that are God's'' open for the Fathers to fill from Genesis.

The other two interpretations answer the same question from other clauses of
the scene. For the second, ``the things that are God's'' are first what Caesar
may not claim: whatever he asks that belongs to God, St John Chrysostom says,
is ``no longer Caesar's tribute, but the devil's'' (\work{Homilies on Matthew}
70.2), so that God's due marks the boundary of Caesar's. For the third, what God
is owed includes the heart of the one who pays, and a correct payment from a
divided heart falls short. The psalm gives the first and third an image they
share. On the summons to bring the Lord glory, Augustine says ``Ascribe them not
unto yourselves'', and on his courts: ``Enter with an humble heart into the
house of God, and thou hast entered with a present'' (\work{Enarrationes in
Psalmos} 95.8--9).

\subsection{Render, not give: the debt to Caesar and where it ends}

The questioners asked whether it was lawful to \emph{give} tribute; Jesus told
them to \emph{render} it. Chrysostom builds on the change of verb: ``For this is
not to give but to render, and this He shows both by the image, and by the
superscription.'' The second clause guards the first: ``Then that they might
not say, Thou art subjecting us to men, He added, `And unto God the things that
are God's.' For it is possible both to fulfill to men their claims and to give
unto God the things that are due to God from us'' (70.2). He cites Romans 13:7,
from a chapter this Mass does not read, and on that chapter says that ``this
subjection is not subversive of religion'' (\work{Homilies on Romans} 23). St
Justin Martyr had told the emperors that Christians pay their taxes ``as we have
been taught by Him'', rendering worship ``to God alone'' (\work{First Apology}
17). Aquinas gives the double belonging its shortest form, ``as if he said: You
are God's, and Caesar's'', and adds that what we have from the world binds us
to keep peace with it, as Baruch bids the exiles pray for the king of Babylon
(\work{Super Matthaeum} c.~22).

The same sovereignty sets the limit. Chrysostom says so in the next breath:
Jesus ``is speaking only of those things, which are no detriment to
godliness''; beyond that line the claim ``is no longer Caesar's tribute, but
the devil's''. Augustine gives the limit a case. ``Julian was an infidel
Emperor, an apostate, a wicked man, an idolater; Christian soldiers served an
infidel Emperor; when they came to the cause of Christ, they acknowledged Him
only who was in heaven.'' They refused the incense and obeyed the order to
march: ``They distinguished their everlasting from their temporal master; and
yet they were, for the sake of their everlasting Master, submissive to their
temporal master'' (\work{Enarrationes in Psalmos} 124.7). Their obedience and
their refusal had one source. On the second interpretation Caesar's claim is
real because it lies within a world where the Lord alone is God, and for the
same reason it is bounded. The Fathers' cases are tribute, incense and the
march of soldiers; they teach the order of two masters, not a form of
government.

Not every witness grounds the debt in the same way, and here the first and
second interpretations draw on the same passage of Ambrose for different ends.
Hilary makes the obligation depend on holding what is Caesar's: \latin{Si enim
nihil eius penes nos resederit, conditione reddendi ei quae sua sunt non
tenebimur}, if nothing of his remains with us, we are not held to render him
his own (23.2). Ambrose draws the consequence for the disciple: ``If you wish to
owe nothing to an earthly king, leave all that is yours and follow Christ''
(IX.35). For them the debt follows possession. Chrysostom answers otherwise.
Paul's ``let every soul be subject'' is written ``to show that these
regulations are for all, even for priests, and monks'', and it binds ``if thou
be an Apostle even, or an Evangelist, or a Prophet'' (\work{Homilies on
Romans} 23): the monk and the Apostle, who have done what Ambrose bids, are
named among those who owe. Augustine places the ground in God's ordering:
``thus hath God ordained His Church, that every power ordained in the world may
have honour, and sometimes from those who are better than those in power''
(124.7). He does not take up possession; the honour he describes rests on
God's ordinance and not on what the one who pays it holds, and for him, as for
Chrysostom, its limit is fixed by what belongs to God. All four hold the
tribute lawful and God's claim first. Hilary and Ambrose on one side and
Chrysostom on the other differ on whether one who keeps nothing of the world's
goods still owes Caesar, and so on what the debt rests on. The first
interpretation takes Ambrose for its Christ, the image of God who
owes Caesar nothing of his own; the second keeps him as one side of that
disagreement. The first would be satisfied by a hearer who pays every tax and gives God himself;
the second asks in addition where the tax must be refused, and finds the answer
in a soldier's refusal of incense.

\subsection{A gentile king called the Lord's anointed}

The first reading raises its own difficulty: how can a man be the Lord's
anointed and not know the Lord? Jerome answers from the title. Cyrus is called
\latin{christus}, ``anointed of the Lord, which was the mark of royal power
among the Hebrews; as among us the diadem and the purple are given to emperors
alone'' (\work{Commentarii in Isaiam} XII); Aquinas, ``to Christ, that is, to
the king \ldots\ because of his royal dignity'' (\work{Expositio super
Isaiam} c.~45).
The title names an office, not a faith. Both ask what the Lord wanted of this
king's power. Jerome glosses ``thou hast not known me'' as ``you worshipped
images, not God'', and reads the omitted treasures as given that he ``might
perceive the one God through his benefits''; Aquinas names the fruit required
of Cyrus himself, \latin{suae cognitionis}, the knowledge of God. Chrysostom
adds that God moved Cyrus, ``disposing him to let the Jews go'', without
destroying his free will (\work{Homilies on Matthew} 6).

The verse has a second history. Cyprian's testimonies, Novatian and the
\work{Epistle of Barnabas} read it ``to Christ my Lord''; as the editor of
Novatian's English notes, a copyist had written the Greek \textit{Kyri\=o},
``Lord'', for the name \textit{Kyr\=o}.
Jerome calls it the error of ``very many, not only of the Latins but of the
Greeks'', and will not let ``thou hast not known me'' be said of the one
``through whom the world was reconciled to God''. The Lectionary's Latin,
\latin{de uncto meo Cyro}, excludes the variant. A narrower difference remains
in v.~4, where the Clementine Vulgate reads \latin{assimilavi te}, ``I have made
a likeness of thee'': Jerome reads it of Cyrus going before Christ ``as a type
and image'', Aquinas of the other saints whose names were foretold, with no
type, and the Lectionary's English follows the Hebrew and has no likeness.

The three interpretations give this oracle three weights. For the second it is
the argument: the one relation between the readings that the liturgical books
declare, a God who uses the power of those who do not know him, set beside the
coin of a ruler who does not know him either. Cyrus is the Lord's instrument,
not his type, and the join of the two passages is the \work{Ordo}'s pairing and
this interpretation's inference from it, for the Fathers speak to each passage
separately. For the first the oracle contributes least: a likeness and a name
given by God would frame the coin's image and inscription, but only on the
Vulgate's word and Jerome's reading of it. For the third it contributes by
contrast. Cyrus does not know the Lord and serves his purpose; the questioners
say ``we know'' and tempt. Knowing God appears in what is done, not in what is
said, and Cyrus is no model of sincerity.

The psalm carries ``none besides me'' to the nations, and the second
interpretation hears it as the Lord's reign. The gods of the gentiles, St Robert
Bellarmine says, ``are not true gods, but demons''; God ``began to reign,
through faith, among the gentiles, from the coming of the Messias, where the
devil previously reigned''. Augustine's Old Latin psalter read ``the Lord hath
reigned from the wood'', which the appointed Latin does not, and on his own
text he says that ``Christ already possesseth regions where the Roman empire
hath never yet reached'' (95.2). Bellarmine's Judge will ``have no regard of
persons'', the quality the questioners had praised in Jesus. The first
interpretation hears in the psalm the gift of a humble heart; the third, the
same heart judged by the Lord who judges with justice, and Augustine's name
for its contrary: ``What is the whited wall, save hypocrisy, that is,
pretence? Without it is bright, within it is dirt'' (95.3).

\subsection{True words from a divided heart, and the one who knows it}

The third interpretation begins from the Gospel's own sequence. The praise that
opens the question is accurate, and two verses later Jesus calls the speakers
hypocrites. If the charge does not concern the truth of what they said, what
does it concern? Four witnesses answer with four contrasts. Jerome: the
hypocrite \latin{aliud opere agit, et aliud voce praetendit}, does one thing in
deed and puts forward another with his voice. Chrysostom: they testify the
truth ``not with an upright mind, nor willingly'', and ask not what is good or
lawful but ``What thinkest Thou?'' (70.1). Aquinas: ``evil men begin from
flattery''; the praise is true because Christ is ``joined to God by union'';
and \latin{hypocritae sunt proprie qui aliud in ore habent, aliud in
corde}, ``hypocrites, properly speaking, are those who have one thing in the
mouth, another in the heart''. Augustine: ``because in one way ye speak, in another
way judge, hypocrites ye are'' (\work{Enarrationes in Psalmos} 57.8). Deed
against voice, mind against testimony, mouth against heart, speech against
judgment: none makes the fault the words' content. The Gospel names those who
came, disciples of the Pharisees with the Herodians, and the vice belongs to
anyone who praises in order to use.

Matthew says that Jesus answered ``knowing their wickedness'' (22:18), and the
witnesses ask what kind of knowing that was. Hilary: Christ asked for the denarius
``knowing the inner secrets of their thoughts, for God sees everything hidden
within men'' (23.1). Aquinas: he answered first their mind, then their words, ``because it belongs
to man to answer words, but to God to answer the mind''.
Here the interpretation meets its difficulty, for Jerome states the same
discernment as a rule for every answerer: ``the first virtue of one who answers
is to know the mind of those who ask''. If that is all, the scene teaches prudence
and proves nothing of Christ. Chrysostom answers from the scene. Their words
gave nothing away, ``full of much respect \ldots\ but being God, He was
deceived by none of these things. Wherefore they also ought to have conjectured,
that the rebuke was not the result of conjecture, but a sign of His knowing
their secret thoughts'' (70.1); and so ``He gave them proof of His Godhead, by
revealing the secrets of their hearts'' (70.2). The christological claim rests
on Chrysostom, Hilary and Aquinas, and on the concealment Chrysostom names.
Jerome's sentence is a rule of prudence, and it belongs with Aquinas's lesson
from the same verse, ``that we ought not to believe flatterers''. The two
answers are complementary: a sign of who Christ is, and a lesson for every
hearer.

On the counter-question the four agree, and all three interpretations use the
agreement. Jesus asks whose image the coin bears, and the answer comes out of
the questioners' own mouths. Augustine: they ``gave a kind of answer to
themselves'', and ``from their answer He drew this inference'', unlike the
question of his authority, where the answer was withheld (\work{De sermone
Domini in monte} II.20.70). Chrysostom: ``it is His will to cause them to be bound by
their own answers''. Jerome: ``so that tempters are refuted above
all by their own words''. Aquinas: ``from the answer he draws out the truth''.
The first
interpretation takes the result, the coin declaring its owner in the
questioners' own words; the second, the verb Jesus draws from their answer,
render and not give; the third, the manner, the hypocrite's own mouth giving
the verdict, so that the divided speaker is answered by the words he used to
hide himself.

The remedy stands in the acclamation's own sentence. Paul wants the
Philippians ``blameless and sincere children of God'', \latin{simplices},
before they shine as lights, and Aquinas comments, \latin{simplices simus sicut
filii Dei; quod est, quando intentio est ad unum}, ``let us be simple as sons
of God, which is when the intention is toward one thing'', with James's ``double
minded man'' (Jas 1:8; \work{Super ad Philippenses} c.~2, lect.~4). The word stands in
an earlier clause of v.~15, outside the acclamation's words; but the light Paul
describes is the light of children whose intention is single.

\subsection{Will, majesty and heart: the Missal's prayers in three lights}

The Missal's antiphons and orations serve this Sunday in all three years of the
cycle, and each interpretation hears them through the words they share with
the Gospel. The Collect asks for a will devoted to God \latin{semper} and for
service of his majesty \latin{sincero corde}. For the first it asks for the will
Hilary names among the things owed to God; for the second the word stands
beside the Gospel's Caesar, since \latin{maiestas} belonged to gods and rulers
alike, ``of the gods'' and ``as a title of honor of the Roman emperors'' (Lewis
and Short, s.v.\ \latin{majestas}), and the majesty served first is God's,
with every other service, Caesar's included, ordered under it; for the third \latin{sincero corde} is the precise contrary of what the
Gospel exposes, and \latin{semper} asks for devotion that does not wait on
usefulness.

The Prayer over the Offerings calls the gifts God's before they are offered.
For the first, Jerome's oblations and Aquinas's bread and wine are the literal
counterpart of the self rendered; for the second, the free mind is that of
those who pay Caesar without belonging to him; for the third, a mind free of
self-interest, the opposite of the questioners' respect put to use. The Prayer
after Communion asks to be aided by temporal benefits and instructed by eternal
ones. The first puts the coin with the temporal and the self with the eternal;
the second finds there the order of goods in which Caesar's coin helps and does
not teach, the order Aquinas's ``You are God's, and Caesar's'' names; the third
hears \latin{erudiamur} as a formation of the heart, which men asking ``What
thinkest Thou?'' did not seek.

The Entrance Antiphon is a prayer of trust. Augustine says that the apple of
the eye, small as it is, is that ``whereby the light is distinguished from the
darkness'', and that the wings are ``the defence of Thy love and mercy''
(\work{Enarrationes in Psalmos} 16.8). For the first interpretation the person
who is God's coin asks to be guarded as what is most God's own; for the second
the protection is God's and not a ruler's, and the unsung opening of v.~8 asks
it against those who resist God's right hand, while in the first reading the
Lord grasps the right hand of a king; for the third the psalm's first verse,
unsung, prays from lips without deceit, ``that I may not in judging utter aught
else than I understand in Thee'' (16.2).

\subsection{A church among the nations: the Apostle, the acclamation and the Communion}

The second reading and the acclamation follow courses of their own, and each
interpretation reads them beside the Gospel. For the first, the church ``in
God'' describes a walk, ``If any one walks not according to God, he cannot be
said to be `in God' '' (Chrysostom, \work{Homilies on First Thessalonians} 1),
and its faith, charity and hope are the self rendered in deeds. For the second,
the same homily recalls Jason dragged ``before the rulers of the city'', on the
charge, in Acts, of saying ``that there is another king, Jesus'' (17:7). For the
third, Chrysostom says that God permitted persecution from the first ``that no
one might say, that the Preaching was established lightly or by flattery'', and Paul's preaching
``was not in loftiness of speech, but in power'' (Aquinas, \work{Super I ad
Thessalonicenses} c.~1, lect.~1). The acclamation's lights are likewise
glossed three ways: for the first, a borrowed light, since the faithful are
light ``inasmuch as they have something of the light of him who was the
light of men'' (Aquinas), as an image is a borrowed likeness; for the second,
stars that shine
the brighter in the night, ``whilst thou holdest straight in the midst of the
crooked'' (Chrysostom, \work{Homilies on Philippians} 8); for the third, the
single intention above.

The Communion antiphons are alternatives, and what rests on one holds only
where it is sung. The first names the eyes of the Lord upon those who fear him:
for the first interpretation, the God who, Augustine says, ``looks for His in
man'' (\work{Tractatus} 41.2); for the
second, the God of a psalm whose unsung v.~16 says ``The king is not saved by a
great army'', on which Aquinas writes that no temporal power can bring men to the
salvation of the just (\work{Postilla super Psalmos} 32); for the third, the
God who sees, and Bellarmine's ``Fear, combined with
hope, is the mark of real love'', where the questioners feigned respect. The
second antiphon names the Son of Man who gives his life as a ransom. For the
first interpretation he is Ambrose's Christ, rendering himself entirely, who,
St Bede says, took the form of a servant \latin{ut pro mundo sanguinem
funderet}, to pour out his blood for the world; for the
second, the antiphon's unsung context, the rulers of the gentiles who lord it
over them, makes him the measure of every authority, and Bede reads the passage
as teaching that the height of the virtues comes ``not by power, but by
humility'', while Chrysostom gives the King's own reason: ``even my life did I
give a ransom; and for whom? For enemies'' (\work{Homilies on Matthew} 65.4).
For the third, the truthful teacher whom the questioners praised in order to
trap gives his life for many.

\subsection{Three answers, four senses}

The three share their literal ground, the whole scene from the plot to the
verdict, and differ in the words each makes decisive: ``Whose image and
inscription is this?'' for the first, ``Render \ldots\ to Caesar \ldots\ and to
God'' with the oracle to Cyrus for the second, ``Why do you tempt me, ye
hypocrites?'' with the praise before it for the third. Their allegorical senses
find three aspects of one Christ: the image of God who bears no image of Caesar
and restamps the worn coin; the King who reigns among the gentiles by faith
and, in Augustine's psalter, ``from the wood'', with Cyrus a figure of him only
in Jerome; and the Truth who reads hearts as only God does, with a Church that
shines by borrowed light. Where the second Communion antiphon is sung, all
three meet in the Son of Man who gives his life as a ransom.

Their moral senses ask three things of one hearer: to render himself, the will
devoted and the heart undivided; to render civil dues as justice requires and
refuse whatever is against God, as the soldiers under Julian did; to be single,
giving and believing no flattery. The first and third can be preached together
without strain: render yourself, with an undivided heart. The second adds a
public dimension the others leave implicit, and with it the Fathers'
disagreement about what the debt to Caesar rests on, which the others need not
settle. Their anagogical senses look to one end from three sides: the image
restored and seen, when, as Augustine joins the saying to 1~John 3:2, ``we
shall see Him as He is'' and ``then shall the truth be reproduced in us''; the
Lord who ``will judge the people with justice'' (Ps 95:10), when the rod of the
wicked over the just, felt ``for a season'', ``will not be there for ever''
(Augustine, \work{Enarrationes in Psalmos} 124.8); and the Judge of the world
from whom, Hilary says, nothing hidden within men is hidden.

Each carries its own difficulty. The first rests on an inference the saying
does not state; the second on a pairing the liturgical books declare and no
Father makes; the third must meet one of its own witnesses, Jerome, who treats
the reading of hearts as a rule for any respondent. Together they fill the
Gospel's unexplained clause in three ways that do not exclude one another: with
the self that bears God's image, with a debt to Caesar that God's claim both
grounds and bounds, and with a heart that means what it says.
